\section{Safe head exchange and an unsafe modal shortcut}
\label{sec:head-exchange}

\subsection{A uniform exchange translation}

Write $I(A,B)=4\pi(A,B)$ for numerical implication. The retained
dominated $\rho_8$ basis derives the pure implicational theorem
\[
 X(A,B,C)=I\bigl(I(A,I(B,C)),I(B,I(A,C))\bigr)
\]
uniformly in all numerical $A,B,C$. The source is
\texttt{PA196/Logic/ImplicationBasis.lean}, whose
\texttt{Dominated.exchange} is obtained from
\texttt{M6.Logic.RestrictedImplication.exchange}. Its abstract rule
interface contains only modus ponens, restricted $\rho_8$, reflexivity
of the numerical domination relation, and its left/right implication
bound constructors. It uses none of the first six nonpropositional heads.

Consequently any selection of heads of the shape $I(A,I(B,C))$ may be
replaced by $I(B,I(A,C))$. For a forward proof translation, replace an
old head row by the new head, a derivation of $X(B,A,C)$, and one modus
ponens step. For the reverse translation use the old head,
$X(A,B,C)$, and one modus ponens step. Translate unchanged head rows
literally and relocate antecedent indices on MP and generalization rows.
A bounded scan through the original proof list constructs the translated
list; correctness is quantifier-free induction on the scanned prefix.
The translations preserve the last numerical formula exactly, including
the designated contradiction code~1.

This applies separately and jointly to membership congruence heads 3
and~4 and distribution head~6. Since every formula code is preserved,
the special alias $I(0,0)=0$ causes no difficulty. The argument neither
uses a syntax recoding nor assumes disjoint formula constructors at zero.

\subsection{Why the superficially equivalent dual of head 1 is unsafe}

The old first head is the modal Brouwer schema
\[
 (\exists_D\forall_D H)\to H,
\qquad N(A)=I(A,1),\qquad
 \exists_D A=N(U(D,N(A))).
\]
Replacing it with $U(D,H)\to H$, or with its familiar modal dual
\[
 H\to U(D,N(U(D,N(H)))),
\]
does not preserve the old calculus by the usual textbook argument.
That argument requires double-negation elimination relative to the
designated numerical formula~1. In the original numerical calculus,
explosion relative to~1 is itself derived using the old first head.
This dependency is explicit in
\texttt{PA196/Logic/Deduction.lean},
\texttt{Context.explosion} and
\texttt{doubleNegationElimination}.

Here is a countervaluation to both proposed derivability transfers.
Use the Boolean truth values, interpret $I$ by ordinary implication,
interpret every $U(a,B)$ by the value of $B$, assign code~0 truth,
assign $M(0,0)$ truth and every other membership atom falsity, and
assign every equality atom falsity. This is well defined by recursion
on formula codes; the zero alias is compatible with Boolean implication.
Because code~1 is true, every $N(A)$ and every defined existential is
true. All retained congruence and witness heads are true: congruence has
a false equality antecedent, and the witness is a defined existential.
Commutation, distribution and the diagonal head are true because $U$
acts as the identity. Every gate is true because it is a numerical
negation. The dominated propositional head is a pure implicational
tautology. Modus ponens and theorem generalization preserve truth.
Both proposed replacement heads are true under this valuation.

However, for $H=M(1,0)$, the old first head has a true antecedent and
false consequent. So does $1\to H$. Thus neither the old first head
nor explosion is derivable from either replacement calculus.
This is a precise obstruction to the proposed proof translation;
it does not by itself settle the consistency strength of the altered
calculus. Retaining the original first head while exchanging the
antecedents of heads 3, 4 or 6 avoids this problem.

